\chapter{Výběr hardwarových komponent}

Při volbě součástkové základny pro navrhovaný palubní počítač byl kladen důraz na spolehlivost v automobilovém prostředí, 
dostatečný výpočetní výkon pro zpracování přerušení v reálném čase, nízkou pořizovací cenu a dobrou dostupnost na trhu. 
Následující podkapitoly shrnují klíčové komponenty zvolené pro realizaci funkčního vzorku.

\section{Mikrokontrolér}

Jako řídicí prvek celého systému byl zvolen modul ESP32 od společnosti Espressif Systems. 
Srdcem tohoto modulu je dvoujádrový 32bitový mikroprocesor Xtensa LX6 s volitelnou taktovací frekvencí až 240 MHz a interní pamětí SRAM o velikosti 520 kB. 
Oproti tradičním osmibitovým mikrokontrolérům (např. ATmega328P z platformy Arduino) nabízí ESP32 zásadní výhody pro aplikaci palubního počítače:

\begin{enumerate}
    \item \textbf{Dvoujádrová architektura a podpora FreeRTOS:}Umožňuje striktně oddělit časově kritické úlohy. Například oddělení měření šířky vstřikovacích impulzů s 
    mikrosekundovým rozlišením a čtení pulzů rychlosti pomocí hardwarových přerušení či periferie Pulse Counter,
    od úloh s nižší prioritou, jako je překreslování grafického displeje a přepočet průměrných hodnot.
    \item \textbf{Hardwarové periferie a konektivita:} Modul disponuje hardwarovou podporou sběrnic I2C, SPI a UART. 
    Integrovaný bezdrátový modul Wi-Fi a Bluetooth Low Energy navíc otevírá možnost budoucího rozšíření systému o bezdrátový export jízdních logů do mobilního 
    telefonu či aktualizaci firmwaru bez nutnosti demontáže přístroje z palubní desky.
\end{enumerate} 

Přestože mikrokontrolér ESP32 obsahuje integrovaný 12bitový analogově-digitální převodník, 
je tato interní periferie známá svou nelinearitou na okrajích měřicího rozsahu a vyšší citlivostí na šum napájecí větve. 
Z tohoto důvodu bylo pro měření analogových veličin rozhodnuto o nasazení externího 16bitového sigma-delta A/D převodníku ADS1115 komunikujícího po sběrnici I2C. 
Převodník ADS1115 disponuje integrovaným zdrojem referenčního napětí, programovatelným zesilovačem a čtyřmi multiplexovanými vstupy, 
což zajišťuje vysokou přesnost a teplotní stabilitu měření.

\section{Zobrazovací jednotka}

Pro přehledné zobrazení provozních parametrů za jízdy je vyžadován displej s vysokým kontrastem, širokými pozorovacími úhly a dobrou čitelností jak na přímém slunci, 
tak při noční jízdě.

Pro toto rozhraní byl zvolen grafický monochromatický OLED displej s řadičem SH1106 a rozlišením 128 × 64 obrazových bodů. 
Technologie OLED nevyžaduje celoplošné podsvícení, neboť světlo emitují přímo jednotlivé aktivní pixely. 
Díky tomu displej dosahuje řádově vyššího kontrastu než dobové pasivní LCD displeje a neoslňuje řidiče při nočním provozu. 
Komunikace po dvouvodičové sběrnici I2C navíc minimalizuje počet vodičů vedených k displeji a umožňuje sdílet sběrnici se zmíněným A/D převodníkem ADS1115.

\section{Senzory a napájecí moduly}

Zbytek hardwarové architektury tvoří obvody pro úpravu signálů z vozidla, pomocné senzory a dvoustupňová napájecí kaskáda:

\begin{enumerate}
    \item \textbf{Napájecí část:} Palubní síť vozidla Škoda Felicia vykazuje za běhu alternátoru napětí v rozmezí 13,5 až 14,4 V. 
Pokud by byl pro snížení tohoto napětí na úroveň 3,3 V použit pouze lineární stabilizátor, mohlo by docházet při odběru mikrokontroléru ESP32 k 
tepelným ztrátám v řádu nižších jednotek wattů. Napájecí obvod je proto navržen jako dvoustupňový. První stupeň tvoří spínaný snižující měnič LM2596S-5 s vysokou účinností, 
který sráží palubní napětí na mezilehlou úroveň 5 V. Druhý stupeň tvoří lineární LDO stabilizátor LD1117V33, 
jenž z 5V větve generuje vyhlazené napětí 3,3 V s minimálním zvlněním pro napájení mikrokontroléru, převodníku ADS1115 a displeje.
    \item \textbf{Snímání vstřikovačů a otáček:} Pro galvanické oddělení spínaných impulzů vstřikovacího ventilu od logické části mikrokontroléru byl zvolen 
    fototranzistorový optočlen PC817. Vstupní infračervená LED optočlenu je chráněna omezovacím rezistorem a antiparalelní diodou 1N4001 proti záporným indukčním 
    špičkám vznikajícím při rozepnutí vinutí vstřikovače.
    \item \textbf{Snímání rychlosti:} Pro odečet otáček mechanického náhonu rychloměru byl vybrán unipolární Hallův spínač TLE4905L od společnosti Infineon, 
    vyvinutý pro automobilové aplikace. Snímač reaguje na magnetické pole rotujícího magnetu v přístrojovém štítu a na svém výstupu s otevřeným kolektorem generuje 
    obdélníkové impulzy odpovídající ujeté dráze.
    \item \textbf{Snímání teploty:} Měření venkovní teploty zajišťuje NTC termistor zapojený v odporovém děliči. 
\end{enumerate}


\chapter{Návrh hardwaru}

\chapter{Softwarová implementace}

\chapter{Konstrukční řešení a implementace do vozidla}

\chapter{Oživení, testování a kalibrace}









% \section{Programové řešení}
% Lorem ipsum dolor sit amet, consectetuer adipiscing elit. Nulla pulvinar eleifend sem. Integer in sapien. Etiam sapien elit, consequat eget, tristique non, venenatis quis, ante. In laoreet, magna id viverra tincidunt, sem odio bibendum justo, vel imperdiet sapien wisi sed libero. Phasellus enim erat, vestibulum vel, aliquam a, posuere eu, velit. Aliquam erat volutpat. Nullam faucibus mi quis velit \cite{sr72/2017}.

% \section{Výsledky měření}
% Fusce tellus odio, dapibus id fermentum quis, suscipit id erat. Fusce tellus. Morbi scelerisque luctus velit. In laoreet, magna id viverra tincidunt, sem odio bibendum justo, vel imperdiet sapien wisi sed libero. Quisque porta. Fusce suscipit libero eget elit. Nulla non lectus sed nisl molestie malesuada. Phasellus faucibus molestie nisl. Integer vulputate sem a nibh rutrum consequat. Proin mattis lacinia justo. Phasellus et lorem id felis nonummy placerat. Etiam ligula pede, sagittis quis, interdum ultricies, scelerisque eu. Cras elementum. Aenean placerat. Donec ipsum massa, ullamcorper in, auctor et, scelerisque sed, est. Aliquam ante. Integer imperdiet lectus quis justo. Vivamus ac leo pretium faucibus. Nullam faucibus mi quis velit.

% \subsection{Etiam quis quam}
% Neque porro quisquam est, qui dolorem ipsum quia dolor sit amet, consectetur, adipisci velit, sed quia non numquam eius modi tempora incidunt ut labore et dolore magnam aliquam quaerat voluptatem. Aliquam erat volutpat. Lorem ipsum dolor sit amet, consectetuer adipiscing elit \cite{sr72/2017,pravidla}. Nunc auctor. Neque porro quisquam est, qui dolorem ipsum quia dolor sit amet, consectetur, adipisci velit, sed quia non numquam eius modi tempora incidunt ut labore et dolore magnam aliquam quaerat voluptatem. Maecenas lorem. Maecenas libero. In laoreet, magna id viverra tincidunt, sem odio bibendum justo, vel imperdiet sapien wisi sed libero. Nullam rhoncus aliquam metus.

% \subsubsection{Integer rutrum orci vestibulum}
% Integer rutrum, orci vestibulum ullamcorper ultricies, lacus quam ultricies odio, vitae placerat pede sem sit amet enim. Ut enim ad minim veniam, quis nostrud exercitation ullamco laboris nisi ut aliquip ex ea commodo consequat. Fusce tellus odio, dapibus id fermentum quis, suscipit id erat. Nullam eget nisl. Nunc auctor. Etiam dui sem, fermentum vitae, sagittis id, malesuada in, quam. Fusce dui leo, imperdiet in, aliquam sit amet, feugiat eu, orci. Curabitur vitae diam non enim vestibulum interdum. Aliquam erat volutpat. Pellentesque sapien. Phasellus enim erat, vestibulum vel, aliquam a, posuere eu, velit.

% \subsubsection{Eger rutrum orci westibulum}
% Fusce dui leo, imperdiet in, aliquam sit amet, feugiat eu, orci. Maecenas aliquet accumsan leo. Aliquam ornare wisi eu metus. Cum sociis natoque penatibus et magnis dis parturient montes, nascetur ridiculus mus. Aliquam erat volutpat. Donec iaculis gravida nulla. Sed elit dui, pellentesque a, faucibus vel, interdum nec, diam. Temporibus autem quibusdam et aut officiis debitis aut rerum necessitatibus saepe eveniet ut et voluptates repudiandae sint et molestiae non recusandae. Nulla non arcu lacinia neque faucibus fringilla. Phasellus enim erat, vestibulum vel, aliquam a, posuere eu, velit. Praesent vitae arcu tempor neque lacinia pretium
% \cite{Walter1999,Svacina1999IEEE,RajmicSysel2002}.

% Aliquam erat volutpat. Quisque porta. Integer imperdiet lectus quis justo. Nullam justo enim, consectetuer nec, ullamcorper ac, vestibulum in, elit. Nullam faucibus mi quis velit. Fusce tellus. Fusce consectetuer risus a nunc. Cras pede libero, dapibus nec, pretium sit amet, tempor quis. Morbi imperdiet, mauris ac auctor dictum, nisl ligula egestas nulla, et sollicitudin sem purus in lacus
% \cite{CSN_ISO_690-2022,CSN_ISO_7144-1997,CSN_ISO_31-11}.